% cox e mullender, semáforos.
%
% processo tem um thread com duas pilhas
%
% pike et al, dormir e acordar
\chapter{Escalonamento}
\label{CH:SCHED}

Qualquer sistema operacional provavelmente executará mais processos do que
o computador possui CPUs, então um plano é necessário para compartilhar o tempo
das CPUs entre os processos. Idealmente, o compartilhamento seria transparente para os processos do usuário. Uma abordagem comum é fornecer a cada processo
a ilusão de que tem sua própria CPU virtual por meio do
\indextext{multiplexação}
dos processos nas CPUs de hardware. Este capítulo explica como o xv6
alcança essa multiplexação.

\section{Multiplexação}

O xv6 multiplexa ao alternar cada CPU de um processo para
outro em duas situações. Primeiro, o mecanismo de
\lstinline{sleep} e
\lstinline{wakeup} do xv6 alterna quando um processo faz uma chamada de sistema
que bloqueia (tem que esperar por um evento), tipicamente em \lstinline{read},
\lstinline{wait} ou
\lstinline{sleep}. Em segundo lugar, o xv6 força periodicamente uma troca para lidar com
processos que computam por longos períodos sem bloquear.
As primeiras são trocas voluntárias; as últimas são chamadas involuntárias.
Essa multiplexação cria a ilusão de que
cada processo tem sua própria CPU.

Implementar a multiplexação apresenta alguns desafios. Primeiro, como alternar de um
processo para outro? 
A ideia básica é salvar e restaurar os registradores da CPU,
embora o fato de que isso não pode ser expresso em C torne a tarefa complicada.
Em segundo lugar, como forçar
as trocas de uma maneira que seja transparente para os processos do usuário? O xv6 usa a
técnica padrão em que os interruptores de um timer de hardware conduzem as trocas de contexto.
Em terceiro lugar, todas as CPUs alternam entre o mesmo conjunto de processos, então um
plano de bloqueio é necessário para evitar competições. Em quarto lugar, a
memória e outros recursos de um processo devem ser liberados quando o processo sai,
mas ele não pode fazer tudo isso por conta própria
porque (por exemplo) não pode liberar sua própria pilha de núcleo enquanto ainda a estiver usando.
Em quinto lugar, cada CPU de uma máquina multi-core deve lembrar qual
processo está executando para que as chamadas de sistema afetem o estado de núcleo do processo correto.
Finalmente, \lstinline{sleep} e \lstinline{wakeup} permitem que um processo desista da
CPU e aguarde ser despertado por outro processo ou interrupção.
É necessário ter cuidado para evitar competições que resultem na
perda de notificações de despertar.

%% 
\section{Código: Troca de Contexto}
%% 

\begin{figure}[t]
\center
\includegraphics[scale=0.5]{fig/switch.pdf}
\caption{Mudança de um processo de usuário para outro. Neste exemplo, o xv6 roda com uma CPU (e, portanto, uma thread de escalonamento).}
\label{fig:switch}
\end{figure}

A Figura~\ref{fig:switch} 
esboça os passos envolvidos na troca de um
processo de usuário para outro:
uma trap (chamada de sistema ou interrupção) do espaço do usuário para a thread de núcleo do processo antigo,
uma troca de contexto para a thread de escalonamento da CPU atual, uma troca
de contexto para a thread de núcleo de um novo processo, e um retorno de trap
para o processo em nível de usuário.
O xv6 possui threads separadas
(registradores e pilhas salvos)
em que executar o escalonador porque
não é seguro para o escalonador executar na
pilha de núcleo de qualquer processo: outra CPU pode
acordar o processo e executá-lo, e seria um desastre
usar a mesma pilha em duas CPUs diferentes.
Há uma thread de escalonamento separada para cada CPU para lidar
com situações em que mais de uma CPU está executando
um processo que deseja desistir da CPU.

Nesta seção, examinaremos a mecânica de troca
entre uma thread de núcleo e uma thread de escalonamento.

Trocar de uma thread para outra envolve salvar os
registradores da CPU da thread antiga e restaurar os registradores
salvos anteriormente da nova thread; o fato de que
o ponteiro da pilha e o contador de programa
são salvos e restaurados significa que a CPU trocará de pilha e
trocará o código que está executando.

A função
\indexcode{swtch}
salva e restaura registradores para uma troca de thread de núcleo.
\lstinline{swtch}
não conhece diretamente as threads; ela apenas salva e
restaura conjuntos de registradores RISC-V, chamados de
\indextext{contextos}.
Quando é hora de um processo desistir da CPU,
a thread de núcleo do processo chama
\lstinline{swtch}
para salvar seu próprio contexto e restaurar o contexto do escalonador.
Cada contexto está contido em uma
\lstinline{struct}
\lstinline{context}
\lineref{kernel/proc.h:/^struct.context/},
que está contida na
\lstinline{struct proc}
de um processo ou na
\lstinline{struct cpu}
de uma CPU.
\lstinline{swtch}
recebe dois argumentos:
\indexcode{struct context}
\lstinline{*old}
e
\lstinline{struct}
\lstinline{context}
\lstinline{*new}.
Ela salva os registradores atuais em
\lstinline{old},
carrega os registradores de 
\lstinline{new},
e retorna.

Vamos acompanhar um processo através de
\lstinline{swtch} 
para o escalonador.
Vimos no Capítulo~\ref{CH:TRAP}
que uma possibilidade ao final de uma interrupção é que 
\indexcode{usertrap}
chame 
\indexcode{yield}.
\lstinline{yield}
por sua vez chama
\indexcode{sched},
que chama
\indexcode{swtch}
para salvar o contexto atual em
\lstinline{p->context}
e trocar para o contexto do escalonador previamente salvo em 
\indexcode{cpu->context}
\lineref{kernel/proc.c:/swtch..p/}.

\lstinline{swtch}
\lineref{kernel/swtch.S:/swtch/}
salva apenas os registradores salvos pelo chamador;
o compilador C gera código no chamador para salvar
os registradores salvos pelo chamador na pilha.
\lstinline{swtch} conhece o deslocamento de cada
campo de registrador em 
\lstinline{struct context}.
Ela não salva o contador de programa.
Em vez disso,
\lstinline{swtch}
salva o
registrador \lstinline{ra},
que armazena o endereço de retorno de onde
\lstinline{swtch}
foi chamado.
Agora
\lstinline{swtch}
restaura os registradores do novo contexto,
que contém os valores dos registradores salvos por um
\lstinline{swtch} anterior.
Quando 
\lstinline{swtch}
retorna, ele retorna para as instruções apontadas pelo
registrador restaurado
\lstinline{ra}, ou seja,
a instrução de onde a nova thread anteriormente
chamou \lstinline{swtch}.
Além disso, ele retorna na pilha da nova thread,
uma vez que é para lá que o ponteiro da pilha restaurado
aponta.

No nosso exemplo, 
\indexcode{sched}
chamou
\indexcode{swtch}
para trocar para
\indexcode{cpu->context},
o contexto do escalonador por CPU.
Esse contexto foi salvo no momento do passado em que
\lstinline{scheduler}
chamou
\lstinline{swtch}
\lineref{kernel/proc.c:/swtch.&c/}
para trocar para o processo que agora está desistindo da CPU.
Quando o
\indexcode{swtch}
que estamos rastreando retorna,
ele retorna não para
\lstinline{sched}
mas para 
\indexcode{scheduler},
com o ponteiro da pilha na pilha do
escalonador da CPU atual.
%% 
\section{Código: Escalonamento}
%% 

A última seção analisou os detalhes de baixo nível de
\indexcode{swtch}; agora vamos considerar 
\lstinline{swtch}
como dado e examinar 
a troca de uma thread de núcleo de um processo
pelo escalonador para outro processo.
O escalonador existe na forma de uma thread especial por CPU, cada uma executando a
função \lstinline{scheduler}.
Essa função é responsável por escolher qual processo executar em seguida.
Um processo
que deseja desistir da CPU deve
adquirir seu próprio bloqueio de processo
\indexcode{p->lock},
liberar qualquer outro bloqueio que esteja segurando,
atualizar seu próprio estado
(\lstinline{p->state}),
e então chamar
\indexcode{sched}.
Você pode ver essa sequência em
\lstinline{yield}
\lineref{kernel/proc.c:/^yield/},
\texttt{sleep}
e
\texttt{exit}.
\lstinline{sched}
verifica novamente alguns desses requisitos
\linerefs{kernel/proc.c:/if..holding/,/running/}
e então verifica uma implicação:
uma vez que um bloqueio está segurado, as interrupções devem ser desabilitadas.
Finalmente,
\indexcode{sched}
chama
\indexcode{swtch}
para salvar o contexto atual em 
\lstinline{p->context}
e trocar para o contexto do escalonador em
\indexcode{cpu->context}.
\lstinline{swtch}
retorna na pilha do escalonador
como se o
\indexcode{scheduler}
tivesse retornado de
\lineref{kernel/proc.c:/swtch.*contex.*contex/}.
O escalonador continua seu
\lstinline{for}
loop, encontra um processo para executar, 
troca para ele, e o ciclo se repete.

Acabamos de ver que o xv6 mantém
\indexcode{p->lock}
durante chamadas para
\lstinline{swtch}:
o chamador de
\indexcode{swtch}
deve já estar segurando o bloqueio, e o controle do bloqueio passa para o
código para o qual foi feita a troca. Essa disposição é incomum: é mais comum que
a thread que adquire um bloqueio também o libere.
A troca de contexto do xv6 deve quebrar essa convenção porque
\indexcode{p->lock}
protege invariantes nos campos
\lstinline{state}
e
\lstinline{context}
do processo que não são verdadeiros enquanto executa em
\lstinline{swtch}.
Por exemplo, se
\lstinline{p->lock}
não fosse segurado durante
\indexcode{swtch},
uma CPU diferente poderia decidir
executar o processo depois que 
\indexcode{yield}
definiu seu estado como
\lstinline{RUNNABLE},
mas antes que 
\lstinline{swtch}
fizesse com que ele parasse de usar sua própria pilha de núcleo.
O resultado seria duas CPUs executando na mesma pilha,
o que causaria caos.
Uma vez que \lstinline{yield} começou a modificar o estado de um processo em execução
para torná-lo
\lstinline{RUNNABLE},
\lstinline{p->lock} deve permanecer segurado até que as invariantes sejam restauradas:
o ponto de liberação correto mais cedo é após
\lstinline{scheduler}
(executando em sua própria pilha)
limpar
\lstinline{c->proc}.
De maneira semelhante, uma vez que 
\lstinline{scheduler}
começa a converter um processo \lstinline{RUNNABLE} para
\lstinline{RUNNING},
o bloqueio não pode ser liberado até que a thread de núcleo do processo
esteja completamente em execução (após o
\lstinline{swtch},
por exemplo, em
\lstinline{yield}).

O único lugar onde uma thread de núcleo desiste de sua CPU é em
\lstinline{sched},
e ela sempre troca para a mesma localização em \lstinline{scheduler}, que
(quase) sempre troca para alguma thread de núcleo que anteriormente chamou
\lstinline{sched}. 
Assim, se alguém imprimisse os números de linha onde o xv6 troca
threads, observaria o seguinte padrão simples:
\lineref{kernel/proc.c:/swtch..c/},
\lineref{kernel/proc.c:/swtch..p/},
\lineref{kernel/proc.c:/swtch..c/},
\lineref{kernel/proc.c:/swtch..p/},
e assim por diante.
Procedimentos que transferem controle intencionalmente entre si por meio da troca de threads
são às vezes chamados de
\indextext{corrotinas}; 
neste exemplo,
\indexcode{sched}
e
\indexcode{scheduler}
são corrotinas uma da outra.

Há um caso em que a chamada do escalonador para
\indexcode{swtch}
não termina em
\indexcode{sched}.
\lstinline{allocproc} define o registrador de contexto \lstinline{ra}
de um novo processo para
\indexcode{forkret}
\lineref{kernel/proc.c:/^forkret/},
de modo que seu primeiro \lstinline{swtch} "retorna" 
para o início dessa função.
\lstinline{forkret}
existe para liberar o 
\indexcode{p->lock};
caso contrário, como o novo processo precisa
retornar ao espaço do usuário como se retornasse de \lstinline{fork},
poderia começar em
\lstinline{usertrapret}.

\lstinline{scheduler}
\lineref{kernel/proc.c:/^scheduler/} 
executa um loop:
encontra um processo para executar, executa-o até que ele desista, repete.
O escalonador
percorre a tabela de processos
procurando um processo executável, um que tenha
\lstinline{p->state} 
\lstinline{==}
\lstinline{RUNNABLE}.
Uma vez que encontra um processo, define a variável do processo atual por CPU
\lstinline{c->proc},
marca o processo como
\lstinline{RUNNING},
e então chama
\indexcode{swtch}
para começar a executá-lo
\linerefs{kernel/proc.c:/Switch.to/,/swtch/}.

%% 
\section{Código: mycpu e myproc}
%% 

O xv6 frequentemente precisa de um ponteiro para a estrutura \lstinline{proc} do processo atual. Em um uniprocessador, poderia haver uma variável global apontando para o \lstinline{proc} atual. 
Isso não funciona em uma máquina multi-core, uma vez que cada CPU executa um processo diferente. A maneira de resolver esse problema é explorar o fato de que cada CPU tem seu próprio conjunto de 
registradores.

Enquanto uma CPU específica está executando no núcleo, o xv6 garante que o registrador \lstinline{tp} da CPU sempre contenha o hartid da CPU. O RISC-V numera suas CPUs, dando a cada uma um 
\indextext{hartid} único. A \indexcode{mycpu} \lineref{kernel/proc.c:/^mycpu/} usa \lstinline{tp} para indexar um array de estruturas \lstinline{cpu} e retornar a correspondente à CPU atual. Uma 
\indexcode{struct cpu} \lineref{kernel/proc.h:/^struct.cpu/} contém um ponteiro para a estrutura \lstinline{proc} do processo que está atualmente em execução nessa CPU (se houver), registradores 
salvos para a thread de escalonamento da CPU e a contagem de spinlocks aninhados necessários para gerenciar a desativação de interrupções.

Garantir que o \lstinline{tp} de uma CPU mantenha o hartid da CPU é um pouco complicado, uma vez que o código do usuário pode modificar \lstinline{tp}. A \lstinline{start} define o registrador 
\lstinline{tp} no início da sequência de inicialização da CPU, enquanto ainda está em modo de máquina \lineref{kernel/start.c:/w_tp/}. A \lstinline{usertrapret} salva \lstinline{tp} na página do 
trampolim, caso o código do usuário o modifique. Finalmente, a \lstinline{uservec} restaura o \lstinline{tp} salvo ao entrar no núcleo a partir do espaço do usuário 
\lineref{kernel/trampoline.S:/make tp hold/}. O compilador garante nunca modificar \lstinline{tp} no código do núcleo. Seria mais conveniente se o xv6 pudesse perguntar ao hardware RISC-V pelo 
hartid atual sempre que necessário, mas o RISC-V permite isso apenas em modo de máquina, não em modo de supervisor.

Os valores de retorno de \lstinline{cpuid} e \lstinline{mycpu} são frágeis: se o timer interromper e causar a thread a desistir e depois retomar a execução em uma CPU diferente, um valor retornado 
anteriormente não será mais correto. Para evitar esse problema, o xv6 requer que os chamadores desativem as interrupções e só as reativem depois de terminarem de usar a \lstinline{struct cpu} 
retornada.

A função \indexcode{myproc} \lineref{kernel/proc.c:/^myproc/} retorna o ponteiro da \lstinline{struct proc} para o processo que está em execução na CPU atual. A \lstinline{myproc} desativa as 
interrupções, invoca \lstinline{mycpu}, busca o ponteiro do processo atual (\lstinline{c->proc}) da \lstinline{struct cpu}, e então habilita as interrupções. O valor de retorno de 
\lstinline{myproc} é seguro para uso mesmo se as interrupções estiverem habilitadas: se uma interrupção de timer mover o processo chamador para uma CPU diferente, seu ponteiro da \lstinline{struct 
proc} permanecerá o mesmo.
%% 
\section{Dormir e Acordar}
\label{sec:sleep}
%% 

O escalonamento e os bloqueios ajudam a ocultar as ações de uma thread
de outra, mas também precisamos de abstrações que ajudem
as threads a interagir intencionalmente.
Por exemplo, o leitor de um pipe no xv6 pode precisar esperar
que um processo escritor produza dados;
a chamada de um pai para \lstinline{wait} pode precisar esperar
que um filho saia; e
um processo lendo do disco precisa esperar
que o hardware do disco termine a leitura.
O núcleo do xv6 usa um mecanismo chamado dormir e acordar
nessas situações (e muitas outras).
Dormir permite que uma thread de núcleo
espere por um evento específico; outra thread pode chamar acordar 
para indicar que as threads esperando por um evento especificado devem retomar.
Dormir e acordar são frequentemente chamados de 
\indextext{coordenação de sequência}
ou 
\indextext{sincronização condicional}
mecanismos.

Dormir e acordar fornecem uma interface de sincronização
relativamente de baixo nível. Para motivar a maneira como funcionam no xv6,
usaremos esses mecanismos para construir uma sincronização de nível superior
chamada \indextext{semaforo}~\cite{dijkstra65} que
coordena produtores e consumidores
(o xv6 não usa semáforos).
Um semáforo mantém uma contagem e fornece duas operações.
A operação "V" (para o produtor) incrementa a contagem.
A operação "P" (para o consumidor) espera até que a contagem seja diferente de zero,
e então a decrementa e retorna.
Se houvesse apenas uma thread produtora e uma thread consumidora,
e elas fossem executadas em CPUs diferentes,
e o compilador não otimizasse de forma agressiva,
essa implementação seria correta:
\begin{lstlisting}[numbers=left,firstnumber=100]
  struct semaphore {
    struct spinlock lock;
    int count;
  };

  void
  V(struct semaphore *s)
  {
     acquire(&s->lock);
     s->count += 1;
     release(&s->lock);
  }

  void
  P(struct semaphore *s)
  {
     while(s->count == 0)
       ;
     acquire(&s->lock);
     s->count -= 1;
     release(&s->lock);
  }
\end{lstlisting}

A implementação acima 
é cara. Se o produtor agir
raramente, o consumidor passará a maior parte
do seu tempo girando no 
\lstinline{while}
loop esperando por uma contagem não zero.
A CPU do consumidor poderia provavelmente encontrar um trabalho mais produtivo do que 
\indextext{espera ocupada}
ao 
\indextext{consultar}
repetidamente 
\lstinline{s->count}.
Evitar a espera ocupada requer
uma maneira para que o consumidor desista da CPU
e retome apenas após
\lstinline{V}
incrementar a contagem.

Aqui está um passo nessa direção, embora, como veremos, não seja suficiente.
Vamos imaginar um par de chamadas, 
\indexcode{sleep}
e
\indexcode{wakeup},
que funcionam da seguinte forma.
\lstinline{sleep(chan)}
espera por um evento designado pelo valor de
\indexcode{chan},
chamado de 
\indextext{canal de espera}.
\lstinline{sleep}
coloca o processo chamador para dormir, liberando a CPU
para outro trabalho.
\lstinline{wakeup(chan)}
acorda todos os processos que estão em chamadas
para \lstinline{sleep} com o mesmo
\lstinline{chan}
(se houver), fazendo com que suas
chamadas de \lstinline{sleep}
retornem.
Se nenhum processo estiver esperando em
\lstinline{chan},
\lstinline{wakeup}
não faz nada.
Podemos mudar a implementação do semáforo para usar
\lstinline{sleep}
e
\lstinline{wakeup} (mudanças destacadas em amarelo):
\begin{lstlisting}[numbers=left,firstnumber=200]
  void
  V(struct semaphore *s)
  {
     acquire(&s->lock);
     s->count += 1;
     (*@\hl{wakeup(s);}@*)
     release(&s->lock);
  }
  
  void
  P(struct semaphore *s)
  {
    while(s->count == 0)    (*@\label{line:test}@*)
      (*@\hl{sleep(s);}@*)  (*@\label{line:sleep}@*)
    acquire(&s->lock);
    s->count -= 1;
    release(&s->lock);
  }
\end{lstlisting}

A \lstinline{P}
agora desiste da CPU em vez de girar, o que é bom.
No entanto, não é simples projetar
\lstinline{sleep}
e 
\lstinline{wakeup}
com essa interface sem sofrer
do que é conhecido como o problema de \indextext{perda de despertador}.
Suponha que
a \lstinline{P}
descubra que
\lstinline{s->count}
\lstinline{==}
\lstinline{0} 
na linha~\ref{line:test}.
Enquanto
\lstinline{P}
está entre as linhas~\ref{line:test} e~\ref{line:sleep},
a \lstinline{V}
é executada em outra CPU:
ela muda
\lstinline{s->count}
para ser não zero e chama
\lstinline{wakeup},
que não encontra processos dormindo e, portanto, não faz nada.
Agora
\lstinline{P}
continua executando na linha~\ref{line:sleep}:
ela chama
\lstinline{sleep}
e vai dormir.
Isso causa um problema:
a \lstinline{P}
está dormindo esperando por uma chamada de \lstinline{V}
que já aconteceu.
A menos que tenhamos sorte e o produtor chame
\lstinline{V} novamente, o consumidor irá esperar
para sempre, mesmo
que a contagem seja não zero.

A raiz desse problema é que a
invariante de que
\lstinline{P}
dorme apenas quando
\lstinline{s->count}
\lstinline{==}
\lstinline{0}
é violada pela 
execução da 
\lstinline{V}
no momento errado.
Uma maneira incorreta de proteger a invariante seria 
mover a aquisição do bloqueio (destacada em amarelo abaixo) em
\lstinline{P}
para que sua verificação da contagem e sua chamada para \lstinline{sleep}
sejam atômicas:
\begin{lstlisting}[numbers=left,firstnumber=300]
  void
  V(struct semaphore *s)
  {
    acquire(&s->lock);
    s->count += 1;
    wakeup(s);
    release(&s->lock);
  }
  
  void
  P(struct semaphore *s)
  {
    (*@\hl{acquire(\&s->lock);}@*)
    while(s->count == 0)    (*@\label{line:test1}@*)
      sleep(s);             (*@\label{line:sleep1}@*)
    s->count -= 1;
    release(&s->lock);
  }
\end{lstlisting}
Pode-se esperar que esta versão da
\lstinline{P}
evite a perda de despertador porque o bloqueio impede
a \lstinline{V}
de executar entre as linhas~\ref{line:test1} e~\ref{line:sleep1}.
Isso acontece, mas também causa um deadlock:
a \lstinline{P}
mantém o bloqueio enquanto dorme,
então a \lstinline{V} ficará bloqueada para sempre esperando pelo bloqueio.

Vamos corrigir o esquema anterior mudando
a interface de \lstinline{sleep}:
o chamador deve passar o \indextext{bloqueio de condição} para
\lstinline{sleep}
para que ele possa liberar o bloqueio após
o processo chamador ser marcado como dormindo e esperando no
canal de espera.
O bloqueio forçará uma
\lstinline{V} concorrente
a esperar até que a \lstinline{P} tenha terminado de se colocar para dormir,
para que a
\lstinline{wakeup}
encontre o consumidor dormindo e o acorde.
Uma vez que o consumidor acorde novamente,
a \indexcode{sleep}
readquire o bloqueio antes de retornar.
Nosso novo esquema correto de dormir/acordar pode ser usado da seguinte forma (mudança
destacada em amarelo):
\begin{lstlisting}[numbers=left,firstnumber=400]
  void
  V(struct semaphore *s)
  {
    acquire(&s->lock);
    s->count += 1;
    wakeup(s);
    release(&s->lock);
  }

  void
  P(struct semaphore *s)
  {
    acquire(&s->lock);
    while(s->count == 0)
       (*@\hl{sleep(s, \&s->lock);}@*)
    s->count -= 1;
    release(&s->lock);
  }
\end{lstlisting}

O fato de que
\lstinline{P}
segura
\lstinline{s->lock}
impede que 
\lstinline{V}
tente acordá-la entre 
a verificação de
\lstinline{s->count}
de 
\lstinline{P}
e sua chamada para
\lstinline{sleep}.
No entanto, 
\lstinline{sleep}
deve liberar
\lstinline{s->lock}
e colocar o processo consumidor para dormir
de uma maneira que seja atômica do ponto de vista de \lstinline{wakeup},
para evitar perdas de despertador.

%% 
\section{Código: Dormir e Acordar}
%% 

As funções \indexcode{sleep} \lineref{kernel/proc.c:/^sleep/} e
\indexcode{wakeup} \lineref{kernel/proc.c:/^wakeup/} do xv6
fornecem a interface usada no último exemplo acima.
A ideia básica é que \lstinline{sleep} marque o processo atual como
\indexcode{SLEEPING}
e então chame \indexcode{sched} para liberar a CPU;
\lstinline{wakeup} procura um processo dormindo no canal de espera dado
e o marca como 
\indexcode{RUNNABLE}.
Os chamadores de
\lstinline{sleep}
e
\lstinline{wakeup}
podem usar qualquer número mutuamente conveniente como o canal.
O xv6 frequentemente usa o endereço
de uma estrutura de dados do núcleo envolvida na espera.

A \lstinline{sleep}
adquire 
\indexcode{p->lock} \lineref{kernel/proc.c:/DOC: sleeplock1/}
e {\it somente então} libera
\lstinline{lk}.
Como veremos, o fato de que \lstinline{sleep}
segura um ou outro desses bloqueios em todos os momentos é o que
impede uma \lstinline{wakeup} concorrente
(que deve adquirir e manter ambos) de agir.
Agora que
\lstinline{sleep}
segura apenas
\lstinline{p->lock},
pode colocar o processo para dormir registrando
o canal de espera,
mudando o estado do processo para \texttt{SLEEPING},
e chamando
\lstinline{sched} \linerefs{kernel/proc.c:/chan.=.chan/,/sched/}.
Em breve ficará claro por que é crítico que
\lstinline{p->lock} não seja liberado (pelo \lstinline{scheduler}) até depois
que o processo esteja marcado como \texttt{SLEEPING}.

Em algum momento, um processo adquirirá o bloqueio de condição,
definirá a condição que o dorminhoco está esperando,
e chamará \lstinline{wakeup(chan)}.
É importante que \lstinline{wakeup} seja chamado
enquanto segura o bloqueio de condição\footnote{%
%
Estritamente falando, é suficiente se
\lstinline{wakeup}
meramente seguir o
\lstinline{acquire}
(ou seja, poderia-se chamar
\lstinline{wakeup}
após o
\lstinline{release}).%
%
}.
A \lstinline{wakeup}
percorre a tabela de processos
\lineref{kernel/proc.c:/^wakeup\(/}.
Ela adquire o
\lstinline{p->lock}
de cada processo que inspeciona.
Quando \lstinline{wakeup} encontra um processo no estado
\indexcode{SLEEPING}
com um
\indexcode{chan}
correspondente,
ela muda o estado desse processo para
\indexcode{RUNNABLE}.
Na próxima vez que o \lstinline{scheduler} rodar, ele verá que o processo está pronto para ser executado.

Por que as regras de bloqueio para 
\lstinline{sleep}
e
\lstinline{wakeup}
garantem que um processo que está indo para dormir
não perderá um despertar concorrente?
O processo que está indo para dormir segura
o bloqueio de condição ou seu próprio
\lstinline{p->lock} 
ou ambos desde {\it antes} de verificar a condição
até {\it depois} de ser marcado como \texttt{SLEEPING}.
O processo que chama \texttt{wakeup} segura \textit{ambos}
os bloqueios no loop de \texttt{wakeup}.
Assim, o despertador ou torna a condição verdadeira antes
que a thread consumidora verifique a condição;
ou o \lstinline{wakeup} do despertador examina a thread dormindo
estritamente depois que ela foi marcada como \texttt{SLEEPING}.
Então, 
\lstinline{wakeup}
verá o processo dormindo e o acordará
(a menos que algo mais o acorde primeiro).

Às vezes, múltiplos processos estão dormindo
no mesmo canal; por exemplo, mais de um processo
lendo de um pipe.
Uma única chamada para 
\lstinline{wakeup}
acordará todos eles.
Um deles será executado primeiro e adquirirá o bloqueio com o qual
\lstinline{sleep}
foi chamado, e (no caso de pipes) lerá qualquer
dado que estiver esperando.
Os outros processos descobrirão que, apesar de terem sido acordados,
não há dados para serem lidos.
Do ponto de vista deles, o despertar foi "espúrio", e
eles devem dormir novamente.
Por essa razão, \lstinline{sleep} é sempre chamado dentro de um loop que
verifica a condição.

Não há problema se duas chamadas de sleep/wakeup acidentalmente
escolhem o mesmo canal: elas verão despertares espúrios,
mas o looping conforme descrito acima tolerará esse problema.
Grande parte do charme de dormir/acordar é que é tanto
leve (não há necessidade de criar estruturas de dados especiais
para atuar como canais de espera) quanto fornece uma camada
de indireção (os chamadores não precisam saber qual processo específico
estão interagindo).
%% 
\section{Código: Pipes}
%% 

Um exemplo mais complexo que usa \lstinline{sleep}
e \lstinline{wakeup} para sincronizar produtores e
consumidores é a implementação de pipes do xv6.
Vimos a interface para pipes no Capítulo~\ref{CH:UNIX}:
bytes escritos em uma extremidade de um pipe são copiados
para um buffer no núcleo e, em seguida, podem ser lidos da
outra extremidade do pipe.
Capítulos futuros examinarão o suporte a descritores de arquivo
em torno de pipes, mas vamos olhar agora as
implementações de 
\indexcode{pipewrite}
e
\indexcode{piperead}.

Cada pipe
é representado por uma 
\indexcode{struct pipe},
que contém
um 
\lstinline{lock}
e um 
\lstinline{data}
buffer.
Os campos
\lstinline{nread}
e
\lstinline{nwrite}
contam o número total de bytes lidos do
e escritos no buffer.
O buffer se envolve:
o próximo byte escrito após
\lstinline{buf[PIPESIZE-1]}
é 
\lstinline{buf[0]}.
As contagens não se envolvem.
Essa convenção permite que a implementação
distinga um buffer cheio 
(\lstinline{nwrite}
\lstinline{==}
\lstinline{nread+PIPESIZE})
de um buffer vazio
(\lstinline{nwrite}
\lstinline{==}
\lstinline{nread}),
mas isso significa que a indexação no buffer
deve usar
\lstinline{buf[nread}
\lstinline{%}
\lstinline{PIPESIZE]}
em vez de apenas
\lstinline{buf[nread]} 
(e de forma semelhante para
\lstinline{nwrite}).

Vamos supor que chamadas para
\lstinline{piperead}
e
\lstinline{pipewrite}
ocorram simultaneamente em duas CPUs diferentes.
A \lstinline{pipewrite} \lineref{kernel/pipe.c:/^pipewrite/}
começa adquirindo o bloqueio do pipe, que
protege as contagens, os dados e suas
invariantes associadas.
A \lstinline{piperead} \lineref{kernel/pipe.c:/^piperead/}
então tenta adquirir o bloqueio também, mas não consegue.
Ela gira em
\lstinline{acquire} \lineref{kernel/spinlock.c:/^acquire/}
esperando pelo bloqueio.
Enquanto
\lstinline{piperead}
espera,
\lstinline{pipewrite}
percorre os bytes que estão sendo escritos
(\lstinline{addr[0..n-1]}),
adicionando cada um ao pipe em sequência
\lineref{kernel/pipe.c:/nwrite\+\+/}.
Durante esse loop, pode acontecer que
o buffer se encha
\lineref{kernel/pipe.c:/DOC: pipewrite-full/}.
Nesse caso, 
a \lstinline{pipewrite}
chama
\lstinline{wakeup}
para alertar quaisquer leitores em espera sobre o fato
de que há dados esperando no buffer
e então dorme em
\lstinline{&pi->nwrite}
para esperar que um leitor retire alguns bytes
do buffer.
A \lstinline{sleep}
libera 
o bloqueio do pipe
como parte de colocar o
processo da \lstinline{pipewrite} para dormir.

A \lstinline{piperead}
agora adquire o bloqueio do pipe e entra em sua seção crítica:
ela descobre que
\lstinline{pi->nread}
\lstinline{!=}
\lstinline{pi->nwrite}
\lineref{kernel/pipe.c:/DOC: pipe-empty/}
(a \lstinline{pipewrite}
foi dormir porque
\lstinline{pi->nwrite}
\lstinline{==}
\lstinline{pi->nread}
\lstinline{+}
\lstinline{PIPESIZE}
\lineref{kernel/pipe.c:/pipewrite-full/}),
então ela cai no 
loop \lstinline{for}, copia os dados do pipe
\lineref{kernel/pipe.c:/DOC: piperead-copy/},
e incrementa 
\lstinline{nread}
pelo número de bytes copiados.
Esse número de bytes agora está disponível para escrita, então
a \lstinline{piperead}
chama
\lstinline{wakeup}
\lineref{kernel/pipe.c:/DOC: piperead-wakeup/}
para acordar quaisquer escritores em espera
antes de retornar.
A \lstinline{wakeup}
encontra um processo dormindo em
\lstinline{&pi->nwrite},
o processo que estava executando
\lstinline{pipewrite}
mas parou quando o buffer se encheu.
Ela marca esse processo como
\indexcode{RUNNABLE}.

O código do pipe usa canais de sono separados para o leitor e o escritor
(\lstinline{pi->nread}
e
\lstinline{pi->nwrite});
isso pode tornar o sistema mais eficiente no improvável
evento de que haja muitos
leitores e escritores esperando pelo mesmo pipe.
O código do pipe dorme dentro de um loop verificando a
condição de sono; se houver múltiplos leitores
ou escritores, todos, exceto o primeiro processo a acordar
verão que a condição ainda é falsa e dormirão novamente.
%% 
\section{Código: Esperar, Sair e Matar}
%% 

As funções \lstinline{sleep} e \lstinline{wakeup}
podem ser usadas para muitos tipos de espera.
Um exemplo interessante, introduzido no Capítulo~\ref{CH:UNIX},
é a interação entre a \indexcode{exit} de um filho
e a \indexcode{wait} de seu pai.
No momento da morte do filho, o pai pode já
estar dormindo em {\tt wait}, ou pode estar fazendo outra coisa;
neste último caso, uma chamada subsequente para {\tt wait} deve
observar a morte do filho, talvez muito tempo depois de chamar {\tt exit}.
A maneira como o xv6 registra a morte do filho até que {\tt wait}
a observe é que {\tt exit} coloca o chamador no estado \indexcode{ZOMBIE},
onde permanece até que o {\tt wait} do pai o note, mude
o estado do filho para {\tt UNUSED}, copie o status de saída do filho,
e retorne o ID do processo do filho ao pai.
Se o pai sair antes do filho, o 
pai entrega o filho ao
processo \lstinline{init}, que chama perpetuamente {\tt wait};
assim, cada filho tem um pai para limpar após ele.
Um desafio é
evitar condições de corrida e deadlock entre
as chamadas simultâneas de \lstinline{wait}
e \lstinline{exit} do pai e do filho,
bem como chamadas simultâneas de
\lstinline{exit} e \lstinline{exit}.

A \lstinline{wait} começa adquirindo
\lstinline{wait_lock} \lineref{kernel/proc.c:/^wait/},
que atua como o bloqueio de condição
que ajuda a garantir que \lstinline{wait} não perca um \lstinline{wakeup}
de um filho que está saindo.
Em seguida, a \lstinline{wait} escaneia a tabela de processos.
Se encontrar um filho no estado \texttt{ZOMBIE},
ela libera os recursos desse filho e
sua estrutura \lstinline{proc}, copia
o status de saída do filho para o endereço fornecido à \lstinline{wait}
(se não for 0),
e retorna o ID do processo do filho.
Se 
a \lstinline{wait}
encontrar filhos, mas nenhum tiver saído,
ela chama
\lstinline{sleep}
para esperar que algum deles saia
\lineref{kernel/proc.c:/DOC: wait-sleep/},
e depois escaneia novamente.
A \lstinline{wait} frequentemente segura dois bloqueios,
\lstinline{wait_lock} e o \lstinline{pp->lock} de algum processo;
a ordem para evitar deadlock é primeiro \lstinline{wait_lock}
e depois \lstinline{pp->lock}.

A \lstinline{exit} \lineref{kernel/proc.c:/^exit/} registra o status de saída,
libera alguns recursos, chama \lstinline{reparent} para dar seus
filhos ao processo \lstinline{init}, acorda o pai caso
ele esteja em \lstinline{wait}, marca o chamador como um zumbi, e
cede permanentemente a CPU. A \lstinline{exit} segura tanto
\lstinline{wait_lock} quanto \lstinline{p->lock} durante essa
sequência.
Ela segura \lstinline{wait_lock} porque 
é o bloqueio de condição
para o \lstinline{wakeup(p->parent)}, impedindo que um pai em
\lstinline{wait} perca o despertar. A \lstinline{exit} deve segurar
\lstinline{p->lock} também, para impedir que um pai em
\lstinline{wait} veja que o filho está no estado
\lstinline{ZOMBIE} antes que o filho tenha finalmente chamado
\lstinline{swtch}. A \lstinline{exit} adquire esses bloqueios na
mesma ordem que a \lstinline{wait} para evitar deadlock.

Pode parecer incorreto que a \lstinline{exit} acorde o pai
antes de definir seu estado como \lstinline{ZOMBIE}, 
mas isso é seguro:
embora
\lstinline{wakeup}
possa fazer o pai rodar,
o loop em
\lstinline{wait}
não pode examinar o filho até que o
\lstinline{p->lock} do filho
seja liberado pelo {\tt scheduler},
então
a \lstinline{wait}
não pode olhar para
o processo que está saindo até muito depois de
a \lstinline{exit}
ter definido seu estado como
\lstinline{ZOMBIE}
\lineref{kernel/proc.c:/state.=.ZOMBIE/}.

Enquanto a
\lstinline{exit} 
permite que um processo termine a si mesmo,
a \lstinline{kill} \lineref{kernel/proc.c:/^kill/} 
permite que um processo solicite que outro termine.
Seria muito complexo para a
\lstinline{kill}
destruir diretamente o processo vítima, uma vez que a vítima
pode estar executando em outra CPU, talvez
no meio de uma sequência sensível de atualizações nas estruturas de dados do núcleo.
Assim,
a \lstinline{kill}
faz muito pouco: apenas define o
\indexcode{p->killed} da vítima
e, se estiver dormindo, a acorda.
Eventualmente, a vítima entrará ou sairá do núcleo,
momento em que o código em
\lstinline{usertrap}
chamará
\lstinline{exit}
se
\lstinline{p->killed}
estiver definido
(verifica chamando
\lstinline{killed}
\lineref{kernel/proc.c:/^killed/}).
Se a vítima estiver executando no espaço do usuário, ela logo entrará
no núcleo fazendo uma chamada de sistema ou porque o timer (ou
algum outro dispositivo) interrompe.

Se o processo vítima estiver em
\lstinline{sleep},
a chamada da \lstinline{kill} para
\lstinline{wakeup}
fará com que a vítima retorne de
\lstinline{sleep}.
Isso é potencialmente perigoso porque 
a condição que está sendo esperada pode não ser verdadeira.
No entanto, as chamadas do xv6 para
\lstinline{sleep}
são sempre envolvidas em um
loop \lstinline{while}
que re-testa a condição após
\lstinline{sleep}
retornar.
Algumas chamadas para
\lstinline{sleep}
também testam
\lstinline{p->killed}
no loop, e abandonam a atividade atual se estiver definido.
Isso é feito apenas quando tal abandono seria correto.
Por exemplo, o código de leitura e escrita de pipes
\lineref{kernel/pipe.c:/killed.pr/} 
retorna se a flag de morte estiver definida; eventualmente, o
código retornará ao trap, que novamente
verificará \lstinline{p->killed} e sairá.

Alguns loops de 
\lstinline{sleep}
do xv6 não verificam
\lstinline{p->killed} 
porque o código está no meio de uma chamada de sistema de múltiplos passos
que deve ser atômica.
O driver virtio
\lineref{kernel/virtio\_disk.c:/sleep.b/} 
é um exemplo: ele não verifica
\lstinline{p->killed}
porque uma operação de disco pode ser uma das várias
escritas que são todas necessárias para que o sistema de arquivos
fique em um estado correto.
Um processo que é morto enquanto aguarda I/O de disco não
sairá até completar a chamada de sistema atual e
\lstinline{usertrap} vê a flag de morte.

\section{Bloqueio de Processos}

O bloqueio associado a cada processo (\lstinline{p->lock}) é o
bloqueio mais complexo do xv6.
Uma maneira simples de pensar sobre \lstinline{p->lock} é
que ele deve ser mantido ao ler ou escrever qualquer um dos seguintes
campos da \lstinline{struct proc}:
\lstinline{p->state},
\lstinline{p->chan},
\lstinline{p->killed},
\lstinline{p->xstate},
e
\lstinline{p->pid}.
Esses campos podem ser usados por outros processos, ou por threads do escalonador
em outras CPUs, então é natural que eles
devam ser protegidos por um bloqueio.

No entanto, a maioria dos usos de \lstinline{p->lock} protege aspectos de nível superior
das estruturas de dados e algoritmos de processo do xv6. Aqui está
o conjunto completo de coisas que \lstinline{p->lock} faz:

\begin{itemize}

\item Juntamente com \lstinline{p->state}, impede condições de corrida na alocação
  de slots \lstinline{proc[]} para novos processos.

\item Oculta um processo de vista enquanto está sendo criado
ou destruído.

\item Impede que o \lstinline{wait} de um pai colete um
processo que definiu seu estado como \lstinline{ZOMBIE} mas não
cedeu ainda a CPU.

\item Impede que o escalonador de outra CPU decida executar
um processo que está cedendo após definir seu estado como \lstinline{RUNNABLE} mas
antes de terminar \lstinline{swtch}.

\item Garante que apenas o escalonador de uma CPU decida executar um
  processo \lstinline{RUNNABLE}.

\item Impede que uma interrupção de timer faça um processo
ceder enquanto está em \lstinline{swtch}.

\item Juntamente com o bloqueio de condição, ajuda a evitar que \lstinline{wakeup}
esqueça um processo que está chamando \lstinline{sleep} mas não terminou
de ceder a CPU.

\item Impede que o processo vítima da \lstinline{kill} saia
e talvez seja realocado entre a verificação de 
\lstinline{kill} do 
\lstinline{p->pid} e a definição de \lstinline{p->killed}.

\item Torna a verificação e a escrita da \lstinline{p->state} atômicas.

\end{itemize}

O campo \lstinline{p->parent} é protegido pelo bloqueio global
\lstinline{wait_lock} em vez de por \lstinline{p->lock}.
Apenas o pai de um processo modifica \lstinline{p->parent}, embora
o campo seja lido tanto pelo próprio processo quanto por outros
processos que buscam seus filhos. O propósito de 
\lstinline{wait_lock} é atuar como o bloqueio de condição quando
\lstinline{wait} dorme esperando que algum filho saia. Um
filho que está saindo segura seja \lstinline{wait_lock} ou \lstinline{p->lock}
até depois que tenha definido seu estado como \lstinline{ZOMBIE}, acordado
seu pai e cedido a CPU. O \lstinline{wait_lock} também
serializa chamadas simultâneas de \lstinline{exit} por um pai e um filho,
de modo que o processo \lstinline{init} (que herda o filho)
é garantido ser acordado de seu \lstinline{wait}.
O \lstinline{wait_lock} é um bloqueio global em vez de um bloqueio por processo
em cada pai, porque, até que um processo o adquira,
não pode saber quem é seu pai.

%% 
\section{Mundo Real}
%% 

O escalonador do xv6 implementa uma política de escalonamento simples, que executa cada processo
em turnos. Essa política é chamada de
\indextext{round robin}.
Sistemas operacionais reais implementam políticas mais sofisticadas que, por exemplo,
permitem que os processos tenham prioridades. A ideia é que um processo de alta prioridade
executável será preferido pelo escalonador em relação a um processo de baixa prioridade que também
está executável. Essas
políticas podem se tornar complexas rapidamente porque há frequentemente objetivos em competição: por
exemplo, o sistema operacional também pode querer garantir justiça e
alta taxa de transferência. Além disso, políticas complexas podem levar a interações não intencionais, como
\indextext{inversão de prioridade}
e 
\indextext{convoys}.
A inversão de prioridade pode ocorrer quando um processo de baixa prioridade e um processo de alta prioridade
ambos usam um determinado bloqueio, que, quando adquirido pelo processo de baixa prioridade, pode impedir o
processo de alta prioridade de progredir. Um longo convoy de processos que estão esperando
pode se formar quando muitos
processos de alta prioridade estão esperando por um processo de baixa prioridade que adquire um
bloqueio compartilhado; uma vez que um convoy se forma, ele pode persistir por muito tempo.
Para evitar esses tipos de problemas, mecanismos adicionais são necessários em
escalonadores sofisticados.

As funções \lstinline{sleep} e \lstinline{wakeup}
são um método de sincronização simples e eficaz,
mas existem muitos outros.
O primeiro desafio em todos eles é
evitar o problema de "perda de despertadores" que vimos no
início do capítulo.
O \lstinline{sleep} do núcleo Unix original
simplesmente desabilitava interrupções,
o que era suficiente porque o Unix rodava em um sistema com um único CPU.
Como o xv6 roda em multiprocessadores,
ele adiciona um bloqueio explícito ao
\lstinline{sleep}.
O \lstinline{msleep} do FreeBSD
adota a mesma abordagem.
O \lstinline{sleep} do Plan 9
usa uma função de callback que é executada com o bloqueio de escalonamento
mantido logo antes de ir para o sono;
a função serve como uma verificação de última hora
da condição de sono, para evitar perdas de despertadores.
O \lstinline{sleep} do núcleo Linux
usa uma fila de processos explícita, chamada de fila de espera, em vez de
um canal de espera; a fila tem seu próprio bloqueio interno.

Escanear todo o conjunto de processos em
\lstinline{wakeup}
é ineficiente. Uma solução melhor é
substituir o
\lstinline{chan}
em ambos
\lstinline{sleep}
e
\lstinline{wakeup}
por uma estrutura de dados que mantenha
uma lista de processos dormindo nessa estrutura,
como a fila de espera do Linux.
O \lstinline{sleep} e
\lstinline{wakeup} do Plan 9
chamam essa estrutura de ponto de encontro.
Muitas bibliotecas de threads se referem à mesma
estrutura como uma variável de condição;
nesse contexto, as operações
\lstinline{sleep}
e
\lstinline{wakeup}
são chamadas de
\lstinline{wait}
e
\lstinline{signal}.
Todos esses mecanismos compartilham o mesmo
sabor: a condição de sono é protegida por
algum tipo de bloqueio que é liberado atômica e durante o sono.

A implementação de
\lstinline{wakeup}
acorda todos os processos que estão esperando em um determinado canal, e pode ser
o caso de que muitos processos estejam esperando por aquele canal específico. O
sistema operacional escalonará todos esses processos e eles competirão para verificar
a condição de sono. Processos que se comportam dessa maneira são às vezes chamados de
\indextext{manada trovejante},
e é melhor evitá-los.
A maioria das variáveis de condição possui duas primitivas para
\lstinline{wakeup}:
\lstinline{signal},
que acorda um processo, e
\lstinline{broadcast},
que acorda todos os processos em espera.

Semáforos são frequentemente usados para sincronização.
A contagem geralmente corresponde a algo como
o número de bytes disponíveis em um buffer de pipe
ou o número de filhos zumbis que um processo possui.
Usar uma contagem explícita como parte da abstração
evita o problema de "perda de despertador":
existe uma contagem explícita do número
de despertadores que ocorreram.
A contagem também evita o despertar espúrio
e os problemas da manada trovejante.

Terminar processos e limpá-los introduz muita complexidade no xv6.
Em a maioria dos sistemas operacionais, é ainda mais complexo, porque, por exemplo, o
processo vítima pode estar profundamente dentro do núcleo dormindo, e desfazer sua
pilha requer cuidado, uma vez que cada função na pilha de chamadas
pode precisar fazer alguma limpeza. Algumas linguagens ajudam fornecendo
um mecanismo de exceção, mas não C.
Além disso, existem outros eventos que podem fazer um processo que está
dormindo ser acordado, mesmo que o evento que está esperando ainda não tenha ocorrido. Por
exemplo, quando um processo Unix está dormindo, outro processo pode enviar um 
\indexcode{signal}
para ele. Nesse caso, o
processo retornará da chamada de sistema interrompida com o valor -1 e com
o código de erro definido como EINTR. A aplicação pode verificar esses valores e
decidir o que fazer. O xv6 não suporta sinais e essa complexidade não surge.

O suporte do xv6 para
\lstinline{kill}
não é totalmente satisfatório: existem loops de sono
que provavelmente deveriam verificar
\lstinline{p->killed}.
Um problema relacionado é que, mesmo para 
loops de \lstinline{sleep}
que verificam
\lstinline{p->killed},
existe uma condição de corrida entre 
\lstinline{sleep}
e
\lstinline{kill};
o último pode definir
\lstinline{p->killed}
e tentar acordar a vítima logo após o loop da vítima
verificar
\lstinline{p->killed}
mas antes de chamar
\lstinline{sleep}.
Se esse problema ocorrer, a vítima não notará a
\lstinline{p->killed}
até que a condição pela qual está esperando ocorra. Isso pode ser bastante mais tarde
ou até nunca
(por exemplo, se a vítima estiver aguardando entrada do console, mas o usuário
não digitar nenhuma entrada).

Um sistema operacional real encontraria estruturas
\lstinline{proc} livres com uma lista de livres explícita
em tempo constante em vez da busca em tempo linear em
\lstinline{allocproc};
o xv6 usa a varredura linear por simplicidade.
%% 
\section{Exercícios}
%% 

\begin{enumerate}

\item Implemente semáforos no xv6 sem usar
\lstinline{sleep} e \lstinline{wakeup} (mas
é permitido usar bloqueios ocupados).
Escolha alguns dos usos de xv6 de sleep e wakeup e
substitua-os por semáforos.
Avalie o resultado.

\item Corrija a condição de corrida mencionada acima entre
\lstinline{kill}
e 
\lstinline{sleep},
para que um
\lstinline{kill}
que ocorre após o loop de sono da vítima verificar
\lstinline{p->killed}
mas antes de chamar
\lstinline{sleep}
resulta na vítima abandonando a chamada de sistema atual.

\item Desenhe um plano para que cada loop de sono verifique 
\lstinline{p->killed}
para que, por exemplo, um processo que está no driver virtio possa retornar rapidamente do loop
se for morto por outro processo.

\item Modifique o xv6 para usar apenas uma troca de contexto ao mudar de
  uma thread de núcleo de um processo para outra, em vez de passar
  pelo thread do escalonador. A thread que cede precisará selecionar
  a próxima thread por conta própria e chamar \texttt{swtch}. Os desafios serão
  evitar que múltiplas CPUs executem a mesma thread acidentalmente;
  acertar o bloqueio; e evitar deadlocks.

\item Modifique o \lstinline{scheduler} do xv6
para usar a instrução
RISC-V
\lstinline{WFI}
(esperar por interrupção)
quando não houver processos executáveis.
Tente garantir que, sempre que houver processos executáveis esperando
para rodar, nenhuma CPU esteja pausando em \texttt{WFI}.

\end{enumerate}
